% Options for packages loaded elsewhere
\PassOptionsToPackage{unicode}{hyperref}
\PassOptionsToPackage{hyphens}{url}
\documentclass[
]{article}
\usepackage{xcolor}
\usepackage{amsmath,amssymb}
\usepackage{cancel}
\setcounter{secnumdepth}{-\maxdimen} % remove section numbering
\usepackage{iftex}
\ifPDFTeX
  \usepackage[T1]{fontenc}
  \usepackage[utf8]{inputenc}
  \usepackage{textcomp} % provide euro and other symbols
\else % if luatex or xetex
  \usepackage{unicode-math} % this also loads fontspec
  \defaultfontfeatures{Scale=MatchLowercase}
  \defaultfontfeatures[\rmfamily]{Ligatures=TeX,Scale=1}
\fi
\usepackage{lmodern}
\ifPDFTeX\else
  % xetex/luatex font selection
\fi
% Use upquote if available, for straight quotes in verbatim environments
\IfFileExists{upquote.sty}{\usepackage{upquote}}{}
\IfFileExists{microtype.sty}{% use microtype if available
  \usepackage[]{microtype}
  \UseMicrotypeSet[protrusion]{basicmath} % disable protrusion for tt fonts
}{}
\makeatletter
\@ifundefined{KOMAClassName}{% if non-KOMA class
  \IfFileExists{parskip.sty}{%
    \usepackage{parskip}
  }{% else
    \setlength{\parindent}{0pt}
    \setlength{\parskip}{6pt plus 2pt minus 1pt}}
}{% if KOMA class
  \KOMAoptions{parskip=half}}
\makeatother
\usepackage{color}
\usepackage{fancyvrb}
\newcommand{\VerbBar}{|}
\newcommand{\VERB}{\Verb[commandchars=\\\{\}]}
\DefineVerbatimEnvironment{Highlighting}{Verbatim}{commandchars=\\\{\}}
% Add ',fontsize=\small' for more characters per line
\newenvironment{Shaded}{}{}
\newcommand{\AlertTok}[1]{\textcolor[rgb]{1.00,0.00,0.00}{\textbf{#1}}}
\newcommand{\AnnotationTok}[1]{\textcolor[rgb]{0.38,0.63,0.69}{\textbf{\textit{#1}}}}
\newcommand{\AttributeTok}[1]{\textcolor[rgb]{0.49,0.56,0.16}{#1}}
\newcommand{\BaseNTok}[1]{\textcolor[rgb]{0.25,0.63,0.44}{#1}}
\newcommand{\BuiltInTok}[1]{\textcolor[rgb]{0.00,0.50,0.00}{#1}}
\newcommand{\CharTok}[1]{\textcolor[rgb]{0.25,0.44,0.63}{#1}}
\newcommand{\CommentTok}[1]{\textcolor[rgb]{0.38,0.63,0.69}{\textit{#1}}}
\newcommand{\CommentVarTok}[1]{\textcolor[rgb]{0.38,0.63,0.69}{\textbf{\textit{#1}}}}
\newcommand{\ConstantTok}[1]{\textcolor[rgb]{0.53,0.00,0.00}{#1}}
\newcommand{\ControlFlowTok}[1]{\textcolor[rgb]{0.00,0.44,0.13}{\textbf{#1}}}
\newcommand{\DataTypeTok}[1]{\textcolor[rgb]{0.56,0.13,0.00}{#1}}
\newcommand{\DecValTok}[1]{\textcolor[rgb]{0.25,0.63,0.44}{#1}}
\newcommand{\DocumentationTok}[1]{\textcolor[rgb]{0.73,0.13,0.13}{\textit{#1}}}
\newcommand{\ErrorTok}[1]{\textcolor[rgb]{1.00,0.00,0.00}{\textbf{#1}}}
\newcommand{\ExtensionTok}[1]{#1}
\newcommand{\FloatTok}[1]{\textcolor[rgb]{0.25,0.63,0.44}{#1}}
\newcommand{\FunctionTok}[1]{\textcolor[rgb]{0.02,0.16,0.49}{#1}}
\newcommand{\ImportTok}[1]{\textcolor[rgb]{0.00,0.50,0.00}{\textbf{#1}}}
\newcommand{\InformationTok}[1]{\textcolor[rgb]{0.38,0.63,0.69}{\textbf{\textit{#1}}}}
\newcommand{\KeywordTok}[1]{\textcolor[rgb]{0.00,0.44,0.13}{\textbf{#1}}}
\newcommand{\NormalTok}[1]{#1}
\newcommand{\OperatorTok}[1]{\textcolor[rgb]{0.40,0.40,0.40}{#1}}
\newcommand{\OtherTok}[1]{\textcolor[rgb]{0.00,0.44,0.13}{#1}}
\newcommand{\PreprocessorTok}[1]{\textcolor[rgb]{0.74,0.48,0.00}{#1}}
\newcommand{\RegionMarkerTok}[1]{#1}
\newcommand{\SpecialCharTok}[1]{\textcolor[rgb]{0.25,0.44,0.63}{#1}}
\newcommand{\SpecialStringTok}[1]{\textcolor[rgb]{0.73,0.40,0.53}{#1}}
\newcommand{\StringTok}[1]{\textcolor[rgb]{0.25,0.44,0.63}{#1}}
\newcommand{\VariableTok}[1]{\textcolor[rgb]{0.10,0.09,0.49}{#1}}
\newcommand{\VerbatimStringTok}[1]{\textcolor[rgb]{0.25,0.44,0.63}{#1}}
\newcommand{\WarningTok}[1]{\textcolor[rgb]{0.38,0.63,0.69}{\textbf{\textit{#1}}}}
\usepackage{longtable,booktabs,array}
\newcounter{none} % for unnumbered tables
\usepackage{calc} % for calculating minipage widths
% Correct order of tables after \paragraph or \subparagraph
\usepackage{etoolbox}
\makeatletter
\patchcmd\longtable{\par}{\if@noskipsec\mbox{}\fi\par}{}{}
\makeatother
% Allow footnotes in longtable head/foot
\IfFileExists{footnotehyper.sty}{\usepackage{footnotehyper}}{\usepackage{footnote}}
\makesavenoteenv{longtable}
\setlength{\emergencystretch}{3em} % prevent overfull lines
\providecommand{\tightlist}{%
  \setlength{\itemsep}{0pt}\setlength{\parskip}{0pt}}
\usepackage{bookmark}
\IfFileExists{xurl.sty}{\usepackage{xurl}}{} % add URL line breaks if available
\urlstyle{same}
\hypersetup{
  hidelinks,
  pdfcreator={LaTeX via pandoc}}

\author{}
\date{}

\begin{document}

Absolutely. This file is \textbf{Digital SAT Domain B3: Data Analysis},
which your notes estimate at roughly \textbf{5--7 of the 44 Math
questions}. It covers ratios/rates/units, percentages, one-variable
statistics, two-variable data and models, probability, statistical
inference, and study design.

I'll teach it from the ground up. The goal is that after this, you can
look at a B3 question and know:

\textbf{what type it is → what formula or method applies → what SAT trap
to avoid → whether Desmos helps.}

\begin{center}\rule{0.5\linewidth}{0.5pt}\end{center}

\section{Part 0: WHAT DATA ANALYSIS IS REALLY
TESTING}\label{part-0-what-data-analysis-is-really-testing}

B3 looks like many unrelated subjects, but the SAT keeps repeating the
same seven jobs:

{\def\LTcaptype{none} % do not increment counter
\begin{longtable}[]{@{}ll@{}}
\toprule\noalign{}
Question type & What you need to recognize \\
\midrule\noalign{}
\endhead
\bottomrule\noalign{}
\endlastfoot
Ratios/rates/units & ``What quantities are being compared?'' \\
Percentages & ``What is the original/base amount?'' \\
Statistics & ``Center or spread?'' \\
Scatterplots/models & ``What does the relationship/model mean?'' \\
Probability & ``What is my denominator?'' \\
Inference & ``What can this sample tell me?'' \\
Study design & ``Can I generalize? Can I claim causation?'' \\
\end{longtable}
}

Your notes identify percent change, unit conversions,
mean/median/spread, conditional probability, margin of error,
line-of-best-fit interpretation, and random sampling versus random
assignment as the most recurring question families.

\begin{center}\rule{0.5\linewidth}{0.5pt}\end{center}

\section{Part 1: RATIOS, RATES, PROPORTIONS, AND
UNITS}\label{part-1-ratios-rates-proportions-and-units}

\subsection{1.1 Ratio}\label{ratio}

A ratio compares two quantities.

If there are:

\[
12\text{ boys},\quad 8\text{ girls},
\]

then:

\[
\text{boys:girls}=12:8=3:2.
\]

This is a \textbf{part-to-part ratio}.

But if the question asks:

\begin{quote}
What fraction of the students are boys?
\end{quote}

Total students:

\[
12+8=20.
\]

Therefore:

\[
\frac{12}{20}=\frac35.
\]

So:

\[
\boxed{\text{boys:girls}=3:2}
\]

but:

\[
\boxed{\text{boys as fraction of whole}=\frac35}.
\]

This distinction is extremely important.

\begin{center}\rule{0.5\linewidth}{0.5pt}\end{center}

\subsection{1.2 Ratio ``parts'' method}\label{ratio-parts-method}

Suppose:

\[
\text{red:blue}=3:5
\]

and there are 64 objects total.

Total ratio parts:

\[
3+5=8.
\]

Each part represents:

\[
64\div8=8.
\]

Red:

\[
3(8)=24.
\]

Blue:

\[
5(8)=40.
\]

Therefore:

\[
\boxed{24\text{ red},40\text{ blue}}.
\]

\subsubsection{SAT shortcut}\label{sat-shortcut}

When given a ratio:

\[
a:b
\]

and a total \(T\):

\[
\boxed{\text{one part}=\frac{T}{a+b}}
\]

then multiply by whichever part you need.

\begin{center}\rule{0.5\linewidth}{0.5pt}\end{center}

\subsection{1.3 Proportions}\label{proportions}

If:

\[
\frac ab=\frac cd,
\]

then:

\[
\boxed{ad=bc}.
\]

Example:

3 pounds cost \$5.40.

7 pounds cost \(x\).

Set corresponding units together:

\[
\frac{3\text{ lb}}{5.40\text{ dollars}}
=
\frac{7\text{ lb}}{x\text{ dollars}}.
\]

Cross multiply:

\[
3x=37.8.
\]

\[
\boxed{x=12.60}.
\]

\subsubsection{Important}\label{important}

Keep the same type of quantity in the same location.

Do not write:

\[
\frac{3}{5.4}=\frac{x}{7}.
\]

That mixes pounds with dollars.

\begin{center}\rule{0.5\linewidth}{0.5pt}\end{center}

\subsection{1.4 Unit rate}\label{unit-rate}

A unit rate means ``per 1.''

240 miles on 8 gallons:

\[
\frac{240}{8}=30.
\]

Therefore:

\[
\boxed{30\text{ miles per gallon}}.
\]

Whenever you see:

\begin{itemize}
\tightlist
\item
  per hour
\item
  per gallon
\item
  per person
\item
  per square mile
\item
  per dollar
\end{itemize}

you are looking at a rate.

\begin{center}\rule{0.5\linewidth}{0.5pt}\end{center}

\subsection{1.5 Direct variation}\label{direct-variation}

Direct variation:

\[
\boxed{y=kx}.
\]

That means the ratio:

\[
\frac yx
\]

is constant.

Example:

When:

\[
x=3,\qquad y=12.
\]

Then:

\[
k=\frac{12}{3}=4.
\]

So:

\[
y=4x.
\]

At:

\[
x=7,
\]

\[
y=28.
\]

Answer:

\[
\boxed{28}.
\]

\subsubsection{Recognition}\label{recognition}

If \(x\) doubles, \(y\) doubles.

If \(x\) triples, \(y\) triples.

\begin{center}\rule{0.5\linewidth}{0.5pt}\end{center}

\subsection{1.6 Inverse variation}\label{inverse-variation}

Inverse variation:

\[
\boxed{y=\frac{k}{x}}
\]

or equivalently:

\[
\boxed{xy=k}.
\]

If \(x\) increases, \(y\) decreases.

Example:

\[
x=4,\quad y=6.
\]

Then:

\[
k=4(6)=24.
\]

If:

\[
x=3,
\]

then:

\[
y=\frac{24}{3}=8.
\]

Answer:

\[
\boxed8.
\]

\subsubsection{Recognition}\label{recognition-1}

Direct:

\[
\boxed{\frac yx=\text{constant}}
\]

Inverse:

\[
\boxed{xy=\text{constant}}.
\]

\begin{center}\rule{0.5\linewidth}{0.5pt}\end{center}

\subsection{1.7 Density}\label{density}

The basic relationship:

\[
\boxed{\text{density}=\frac{\text{mass}}{\text{volume}}}.
\]

Rearrange:

\[
\boxed{\text{mass}=\text{density}\times\text{volume}}
\]

and:

\[
\boxed{\text{volume}=\frac{\text{mass}}{\text{density}}}.
\]

The same idea applies to population density:

\[
\boxed{\text{population density}
=
\frac{\text{population}}{\text{area}}}.
\]

Therefore:

\[
\boxed{\text{population}
=
\text{density}\times\text{area}}.
\]

Your source explicitly groups density, average speed, weighted averages,
and direct/inverse variation under the rate-and-unit toolkit.

\begin{center}\rule{0.5\linewidth}{0.5pt}\end{center}

\subsection{1.8 Unit conversion}\label{unit-conversion}

This is a major SAT skill.

The safest method is \textbf{dimensional analysis}.

Suppose:

\[
60\frac{\text{mi}}{\text{hr}}
\]

must become feet per second.

Know:

\[
1\text{ mi}=5280\text{ ft}
\]

and:

\[
1\text{ hr}=3600\text{ sec}.
\]

Write:

\[
60\frac{\text{mi}}{\text{hr}}
\times
\frac{5280\text{ ft}}{1\text{ mi}}
\times
\frac{1\text{ hr}}{3600\text{ sec}}.
\]

Units cancel:

\[
\frac{\cancel{\text{mi}}}{\cancel{\text{hr}}}
\times
\frac{\text{ft}}{\cancel{\text{mi}}}
\times
\frac{\cancel{\text{hr}}}{\text{sec}}.
\]

Leaving:

\[
\frac{\text{ft}}{\text{sec}}.
\]

Calculate:

\[
60\cdot\frac{5280}{3600}=88.
\]

Therefore:

\[
\boxed{88\text{ ft/s}}.
\]

This cancellation-factor method is the universal approach recommended in
the source.

\begin{center}\rule{0.5\linewidth}{0.5pt}\end{center}

\subsection{1.9 Square units}\label{square-units}

This is one of the biggest traps.

Because:

\[
1\text{ ft}=12\text{ in},
\]

then:

\[
1\text{ ft}^2
=
(12\text{ in})^2
=
144\text{ in}^2.
\]

Not:

\[
12\text{ in}^2.
\]

Example:

\[
24\text{ ft}^2.
\]

Convert:

\[
24(144)=3456.
\]

Therefore:

\[
\boxed{3456\text{ in}^2}.
\]

\begin{center}\rule{0.5\linewidth}{0.5pt}\end{center}

\subsection{1.10 Cubic units}\label{cubic-units}

Same principle.

\[
1\text{ yd}=3\text{ ft}.
\]

Therefore:

\[
1\text{ yd}^3=(3\text{ ft})^3=27\text{ ft}^3.
\]

So:

\[
4\text{ yd}^3=4(27)=108\text{ ft}^3.
\]

Answer:

\[
\boxed{108\text{ ft}^3}.
\]

Remember:

\[
\boxed{\text{length conversion: }k}
\]

\[
\boxed{\text{area conversion: }k^2}
\]

\[
\boxed{\text{volume conversion: }k^3}.
\]

\begin{center}\rule{0.5\linewidth}{0.5pt}\end{center}

\subsection{1.11 Average speed}\label{average-speed}

This formula must be memorized:

\[
\boxed{
\text{average speed}
=
\frac{\text{total distance}}
{\text{total time}}
}
\]

You usually \textbf{cannot average the speeds}.

Example:

100 miles at 50 mph.

Time:

\[
\frac{100}{50}=2\text{ hr}.
\]

Then 100 miles at 25 mph.

Time:

\[
\frac{100}{25}=4\text{ hr}.
\]

Total distance:

\[
200.
\]

Total time:

\[
6.
\]

Average:

\[
\frac{200}{6}
=
33.33\text{ mph}.
\]

Not:

\[
\frac{50+25}{2}=37.5.
\]

The arithmetic average of speeds is generally wrong because different
speeds can be maintained for different lengths of time.

\begin{center}\rule{0.5\linewidth}{0.5pt}\end{center}

\subsection{1.12 Weighted average}\label{weighted-average}

Formula:

\[
\boxed{
\bar x=
\frac{\sum w_ix_i}{\sum w_i}
}
\]

Example:

12 students score 80.

8 students score 90.

You cannot calculate:

\[
\frac{80+90}{2}=85
\]

because the group sizes differ.

Instead:

\[
\frac{12(80)+8(90)}{12+8}.
\]

\[
=
\frac{960+720}{20}
\]

\[
=
\frac{1680}{20}
\]

\[
\boxed{84}.
\]

\begin{center}\rule{0.5\linewidth}{0.5pt}\end{center}

\section{Part 2: PERCENTAGES}\label{part-2-percentages}

This is probably the single highest-yield B3 section.

\subsection{2.1 Percent means ``per 100''}\label{percent-means-per-100}

\[
25\%=0.25=\frac14.
\]

To calculate \(p\%\) of \(x\):

\[
\boxed{\frac p{100}x}.
\]

Example:

18\% of 250:

\[
0.18(250)=45.
\]

\begin{center}\rule{0.5\linewidth}{0.5pt}\end{center}

\subsection{2.2 Percent change}\label{percent-change}

Formula:

\[
\boxed{
\text{percent change}
=
\frac{\text{new}-\text{old}}{\text{old}}\times100\%
}
\]

The denominator is always the \textbf{old/original value}.

Example:

80 to 92.

Change:

\[
92-80=12.
\]

Divide by original:

\[
\frac{12}{80}=0.15.
\]

Therefore:

\[
\boxed{15\%\text{ increase}}.
\]

\subsubsection{Wrong denominator}\label{wrong-denominator}

\[
\frac{12}{92}
\]

is not the percent increase.

\begin{center}\rule{0.5\linewidth}{0.5pt}\end{center}

\subsection{2.3 Multiplier method}\label{multiplier-method}

This is the method I want you using on the SAT.

Increase by \(p\%\):

\[
\boxed{1+\frac p{100}}
\]

Decrease by \(p\%\):

\[
\boxed{1-\frac p{100}}.
\]

Examples:

10\% increase:

\[
\times1.10.
\]

23\% increase:

\[
\times1.23.
\]

15\% decrease:

\[
\times0.85.
\]

40\% decrease:

\[
\times0.60.
\]

Your source explicitly recommends turning percent changes into
multipliers and chaining the multipliers rather than adding percentages.

\begin{center}\rule{0.5\linewidth}{0.5pt}\end{center}

\subsection{2.4 ``Increase by'' vs final
percentage}\label{increase-by-vs-final-percentage}

Suppose \$100 increases by 30\%.

Increase:

\[
30.
\]

New value:

\[
130.
\]

Multiplier:

\[
1.30.
\]

Thus:

\[
\boxed{\text{increase by 30\%}=\text{final is 130\% of original}}.
\]

\begin{center}\rule{0.5\linewidth}{0.5pt}\end{center}

\subsection{2.5 Percent increases above
100\%}\label{percent-increases-above-100}

This causes many mistakes.

Suppose something increases by \textbf{150\%}.

Original = 100\%.

Increase = 150\%.

Final:

\[
100\%+150\%=250\%.
\]

Multiplier:

\[
\boxed{2.5}.
\]

Therefore \$20,000 becomes:

\[
20000(2.5)=50000.
\]

A 150\% increase does \textbf{not} mean multiply by 1.5.

Multiply by:

\[
\boxed{2.5}.
\]

\begin{center}\rule{0.5\linewidth}{0.5pt}\end{center}

\subsection{2.6 Successive percentage
changes}\label{successive-percentage-changes}

Suppose \$100 increases 20\%, then decreases 20\%.

First:

\[
100(1.20)=120.
\]

Then:

\[
120(0.80)=96.
\]

So final:

\[
\$96.
\]

That is a:

\[
\boxed{4\%\text{ decrease}}.
\]

Why don't they cancel?

Because the second 20\% is calculated from \textbf{120}, not 100.

Mathematically:

\[
1.2(0.8)=0.96.
\]

\begin{center}\rule{0.5\linewidth}{0.5pt}\end{center}

\subsection{2.7 Multiple
discounts/markups}\label{multiple-discountsmarkups}

Suppose:

40\% markup followed by 25\% discount.

Multipliers:

\[
1.40
\]

and:

\[
0.75.
\]

Combined:

\[
1.40(0.75)=1.05.
\]

So final price is:

\[
105\%
\]

of original.

Net:

\[
\boxed{5\%\text{ increase}}.
\]

Not:

\[
40\%-25\%=15\%.
\]

\begin{center}\rule{0.5\linewidth}{0.5pt}\end{center}

\subsection{2.8 Reverse percentage}\label{reverse-percentage}

Suppose a shirt costs \$45 \textbf{after a 25\% discount}.

After 25\% off, you retain:

\[
75\%=0.75.
\]

So:

\[
0.75(\text{original})=45.
\]

Therefore:

\[
\text{original}=\frac{45}{0.75}=60.
\]

\[
\boxed{\$60}.
\]

General rule:

\[
\boxed{
\text{original}
=
\frac{\text{final}}
{\text{multiplier}}
}
\]

Never ``add 25\% back.'' The 25\% discount was based on a different
number.

\begin{center}\rule{0.5\linewidth}{0.5pt}\end{center}

\subsection{2.9 Percentage points vs percent
change}\label{percentage-points-vs-percent-change}

Suppose an interest rate goes:

\[
4\%\rightarrow6\%.
\]

Percentage-point increase:

\[
6-4=2.
\]

So:

\[
\boxed{2\text{ percentage points}}.
\]

Relative percent increase:

\[
\frac{6-4}{4}=\frac24=0.5.
\]

So:

\[
\boxed{50\%\text{ increase}}.
\]

These are not the same.

\begin{center}\rule{0.5\linewidth}{0.5pt}\end{center}

\subsection{2.10 Percent of a percent}\label{percent-of-a-percent}

30\% of 40\%:

\[
0.30(0.40)=0.12.
\]

Therefore:

\[
\boxed{12\%}.
\]

\begin{center}\rule{0.5\linewidth}{0.5pt}\end{center}

\subsection{2.11 Percent error}\label{percent-error}

Formula:

\[
\boxed{
\text{percent error}
=
\frac{|\text{estimate}-\text{actual}|}
{\text{actual}}\times100\%
}
\]

Always divide by the \textbf{actual value}.

Example:

Estimate \(45\), actual \(50\).

Error:

\[
|45-50|=5.
\]

Then:

\[
\frac5{50}=0.10.
\]

Therefore:

\[
\boxed{10\%}.
\]

\begin{center}\rule{0.5\linewidth}{0.5pt}\end{center}

\subsection{2.12 Compound growth}\label{compound-growth}

Formula:

\[
\boxed{A=P(1+r)^n}.
\]

Suppose:

\[
P=10000,\qquad r=.03,\qquad n=4.
\]

Then:

\[
A=10000(1.03)^4.
\]

Approximately:

\[
\boxed{11255}.
\]

Do not use:

\[
10000(1+4(.03)).
\]

That is simple growth, not compound growth.

\begin{center}\rule{0.5\linewidth}{0.5pt}\end{center}

\subsection{2.13 Finding a growth rate}\label{finding-a-growth-rate}

Suppose:

\[
5000\rightarrow7000
\]

over 3 years with constant annual percentage growth.

Use:

\[
7000=5000(1+r)^3.
\]

Divide:

\[
1.4=(1+r)^3.
\]

Cube root:

\[
1+r=\sqrt[3]{1.4}.
\]

Approximately:

\[
1+r\approx1.119.
\]

Therefore:

\[
\boxed{r\approx11.9\%}.
\]

\begin{center}\rule{0.5\linewidth}{0.5pt}\end{center}

\subsection{2.14 Simple interest}\label{simple-interest}

Formula:

\[
\boxed{I=Prt}.
\]

Total:

\[
\boxed{A=P(1+rt)}.
\]

Example:

\$1,000 at 5\% simple interest for 2 years:

\[
I=1000(.05)(2)=100.
\]

Total:

\[
\boxed{1100}.
\]

Compound annually:

\[
1000(1.05)^2
=
1102.50.
\]

Different formulas.

\begin{center}\rule{0.5\linewidth}{0.5pt}\end{center}

\section{Part 3: ONE-VARIABLE
STATISTICS}\label{part-3-one-variable-statistics}

\subsection{3.1 Mean}\label{mean}

Formula:

\[
\boxed{\bar x=\frac{\sum x_i}{n}}.
\]

Example:

\[
72,85,90,65,88.
\]

Sum:

\[
400.
\]

Divide by 5:

\[
\boxed{80}.
\]

\begin{center}\rule{0.5\linewidth}{0.5pt}\end{center}

\subsection{3.2 Reverse mean trick}\label{reverse-mean-trick}

If mean and number of observations are known:

\[
\boxed{\text{sum}=(\text{mean})(n)}.
\]

Example:

5 numbers have mean 12.

Total:

\[
5(12)=60.
\]

Four numbers are:

\[
8,10,14,16.
\]

Their sum:

\[
48.
\]

Missing number:

\[
60-48=12.
\]

\[
\boxed{12}.
\]

This comes up constantly.

\begin{center}\rule{0.5\linewidth}{0.5pt}\end{center}

\subsection{3.3 Median}\label{median}

First:

\[
\boxed{\text{SORT THE DATA}}.
\]

Odd number of values:

middle value.

Example:

\[
12,7,9,14,3.
\]

Sort:

\[
3,7,9,12,14.
\]

Median:

\[
\boxed9.
\]

Even number:

average the two middle values.

\[
3,7,9,12.
\]

Middle:

\[
7,9.
\]

Median:

\[
\frac{7+9}{2}=8.
\]

\begin{center}\rule{0.5\linewidth}{0.5pt}\end{center}

\subsection{3.4 Mode}\label{mode}

Mode = most frequent value.

Example:

\[
2,4,4,6,7.
\]

Mode:

\[
\boxed4.
\]

You can have:

one mode,

multiple modes,

or no mode.

\begin{center}\rule{0.5\linewidth}{0.5pt}\end{center}

\subsection{3.5 Range}\label{range}

\[
\boxed{\text{range}=\max-\min}.
\]

For:

\[
2,5,8,13,20
\]

range:

\[
20-2=18.
\]

\begin{center}\rule{0.5\linewidth}{0.5pt}\end{center}

\subsection{3.6 Quartiles and IQR}\label{quartiles-and-iqr}

\[
\boxed{\text{IQR}=Q_3-Q_1}.
\]

IQR measures the spread of the \textbf{middle 50\%}.

Example:

\[
1,2,3,4,5,6,7,8,9.
\]

Median:

\[
5.
\]

Lower half:

\[
1,2,3,4.
\]

So:

\[
Q_1=\frac{2+3}{2}=2.5.
\]

Upper half:

\[
6,7,8,9.
\]

\[
Q_3=\frac{7+8}{2}=7.5.
\]

Thus:

\[
\text{IQR}=7.5-2.5=5.
\]

\[
\boxed5.
\]

The file's procedure is: sort, find median, split the list, find
\(Q_1,Q_3\), subtract.

\begin{center}\rule{0.5\linewidth}{0.5pt}\end{center}

\subsection{3.7 Outlier rule}\label{outlier-rule}

Potential lower outlier:

\[
\boxed{Q_1-1.5(\text{IQR})}
\]

Potential upper outlier:

\[
\boxed{Q_3+1.5(\text{IQR})}.
\]

Anything outside those fences is an outlier under the standard rule.

Example:

\[
Q_1=20,\qquad Q_3=40.
\]

IQR:

\[
20.
\]

Upper fence:

\[
40+1.5(20)
=
70.
\]

A value of 100 is above 70.

Therefore it is an outlier.

\begin{center}\rule{0.5\linewidth}{0.5pt}\end{center}

\subsection{3.8 Mean vs median and
outliers}\label{mean-vs-median-and-outliers}

This is crucial.

The \textbf{mean is sensitive} to extreme values.

The \textbf{median is resistant}.

Example:

\[
2,3,4,5,6.
\]

Mean:

\[
4.
\]

Median:

\[
4.
\]

Replace 6 with 100:

\[
2,3,4,5,100.
\]

Median stays:

\[
4.
\]

Mean becomes:

\[
\frac{114}{5}=22.8.
\]

Huge difference.

So:

\[
\boxed{\text{outliers pull the mean}}
\]

but:

\[
\boxed{\text{median changes little or not at all}}.
\]

Your source emphasizes this exact distinction.

\begin{center}\rule{0.5\linewidth}{0.5pt}\end{center}

\subsection{3.9 Distribution shape}\label{distribution-shape}

\subsection{Symmetric}\label{symmetric}

Approximately mirrored.

Usually:

\[
\boxed{\text{mean}\approx\text{median}}.
\]

\begin{center}\rule{0.5\linewidth}{0.5pt}\end{center}

\subsection{Right-skewed}\label{right-skewed}

Long tail to the right.

Large values pull the mean right.

Therefore:

\[
\boxed{\text{mean}>\text{median}}.
\]

Think income.

A few extremely rich people pull the average upward.

\begin{center}\rule{0.5\linewidth}{0.5pt}\end{center}

\subsection{Left-skewed}\label{left-skewed}

Long tail to the left.

Small values pull mean left.

Therefore:

\[
\boxed{\text{mean}<\text{median}}.
\]

\begin{center}\rule{0.5\linewidth}{0.5pt}\end{center}

\subsection{Bimodal}\label{bimodal}

Two peaks.

The mean can lie in a region where very few observations exist.

So the mean may not represent a ``typical'' observation well.

\begin{center}\rule{0.5\linewidth}{0.5pt}\end{center}

\subsection{3.10 Standard deviation}\label{standard-deviation}

You generally do \textbf{not} calculate SD manually on SAT Math.

Understand the concept:

\[
\boxed{\text{SD measures spread around the mean}}.
\]

Small SD:

values clustered tightly.

Large SD:

values farther away.

Example:

A:

\[
95,100,105.
\]

B:

\[
50,100,150.
\]

Both means:

\[
100.
\]

But B has much larger SD because its values are much farther from 100.

\begin{center}\rule{0.5\linewidth}{0.5pt}\end{center}

\subsection{3.11 Same range does NOT imply same
SD}\label{same-range-does-not-imply-same-sd}

Dataset A may have all values packed near center except extremes.

Dataset B may have many observations spread throughout the range.

Same min/max:

same range.

But:

different SD.

Range uses only \textbf{two observations}.

SD uses every observation.

\begin{center}\rule{0.5\linewidth}{0.5pt}\end{center}

\subsection{3.12 Transforming a data set}\label{transforming-a-data-set}

This is very testable.

Suppose you add \(c\) to \textbf{every value}.

Mean:

\[
+c.
\]

Median:

\[
+c.
\]

But differences between values do not change.

So:

\[
\boxed{\text{SD unchanged}}
\]

\[
\boxed{\text{range unchanged}}
\]

\[
\boxed{\text{IQR unchanged}}.
\]

Example:

Data:

\[
10,15,20.
\]

Add 100:

\[
110,115,120.
\]

The center moves by 100.

Spread remains identical.

\begin{center}\rule{0.5\linewidth}{0.5pt}\end{center}

\subsection{\texorpdfstring{3.13 Multiply every value by
\(c\)}{3.13 Multiply every value by c}}\label{multiply-every-value-by-c}

Then all distances also multiply.

So:

mean multiplied by \(c\),

median multiplied by \(c\),

range multiplied by \(|c|\),

IQR multiplied by \(|c|\),

SD multiplied by \(|c|\).

Your notes summarize these transformation rules directly.

\begin{center}\rule{0.5\linewidth}{0.5pt}\end{center}

\subsection{3.14 Weighted mean from frequency
table}\label{weighted-mean-from-frequency-table}

Suppose:

{\def\LTcaptype{none} % do not increment counter
\begin{longtable}[]{@{}lr@{}}
\toprule\noalign{}
Score & Frequency \\
\midrule\noalign{}
\endhead
\bottomrule\noalign{}
\endlastfoot
70 & 3 \\
80 & 5 \\
90 & 2 \\
\end{longtable}
}

Total people:

\[
3+5+2=10.
\]

Weighted total:

\[
3(70)+5(80)+2(90)
\]

\[
=210+400+180
\]

\[
=790.
\]

Mean:

\[
\frac{790}{10}
=
\boxed{79}.
\]

Do not calculate:

\[
\frac{70+80+90}{3}.
\]

\begin{center}\rule{0.5\linewidth}{0.5pt}\end{center}

\subsection{3.15 Histograms}\label{histograms}

A histogram represents \textbf{numerical data grouped into intervals}.

Important:

bars usually touch.

Bar chart:

categorical groups.

Bars can be separated.

\subsubsection{Finding median interval}\label{finding-median-interval}

Suppose frequencies:

{\def\LTcaptype{none} % do not increment counter
\begin{longtable}[]{@{}lr@{}}
\toprule\noalign{}
Interval & Frequency \\
\midrule\noalign{}
\endhead
\bottomrule\noalign{}
\endlastfoot
0--10 & 4 \\
10--20 & 8 \\
20--30 & 12 \\
30--40 & 6 \\
\end{longtable}
}

Total:

\[
30.
\]

Middle observations are 15th and 16th.

Cumulative:

after first bin = 4.

after second = 12.

after third = 24.

Therefore observations 15 and 16 fall in:

\[
\boxed{20-30}.
\]

\begin{center}\rule{0.5\linewidth}{0.5pt}\end{center}

\subsection{3.16 Box plots}\label{box-plots}

A box plot gives you:

\[
\boxed{\min,Q_1,\text{median},Q_3,\max}
\]

or whiskers to the non-outlier extrema, depending on convention.

The box represents:

\[
Q_1\rightarrow Q_3.
\]

That is the middle:

\[
\boxed{50\%}
\]

of the data.

Use:

median to compare center.

IQR to compare spread.

\begin{center}\rule{0.5\linewidth}{0.5pt}\end{center}

\section{Part 4: TWO-VARIABLE DATA}\label{part-4-two-variable-data}

\subsection{4.1 Scatterplots}\label{scatterplots}

Each point represents:

\[
(x,y).
\]

When inspecting a scatterplot, think:

\textbf{direction, form, strength.}

Direction:

positive,

negative,

none.

Form:

linear,

curved/nonlinear.

Strength:

tight or scattered.

\begin{center}\rule{0.5\linewidth}{0.5pt}\end{center}

\subsection{4.2 Positive correlation}\label{positive-correlation}

As \(x\) increases:

\[
y\text{ tends to increase}.
\]

Example:

study time and test score.

\begin{center}\rule{0.5\linewidth}{0.5pt}\end{center}

\subsection{4.3 Negative correlation}\label{negative-correlation}

As \(x\) increases:

\[
y\text{ tends to decrease}.
\]

Example:

speed for a fixed-distance trip and travel time.

\begin{center}\rule{0.5\linewidth}{0.5pt}\end{center}

\subsection{\texorpdfstring{4.4 Correlation coefficient
\(r\)}{4.4 Correlation coefficient r}}\label{correlation-coefficient-r}

\[
\boxed{-1\le r\le1}.
\]

Closer to:

\[
1
\]

means stronger positive \textbf{linear} association.

Closer to:

\[
-1
\]

means stronger negative linear association.

Closer to:

\[
0
\]

means weak/no linear association.

Important:

\[
r=0
\]

doesn't necessarily mean no relationship whatsoever.

There could be a strong nonlinear pattern.

\begin{center}\rule{0.5\linewidth}{0.5pt}\end{center}

\subsection{4.5 Line of best fit}\label{line-of-best-fit}

General form:

\[
\boxed{y=mx+b}.
\]

Slope:

\[
\boxed{m=\text{predicted change in }y\text{ for each +1 in }x}.
\]

Intercept:

\[
\boxed{b=\text{predicted }y\text{ when }x=0}.
\]

Your notes define slope and intercept this way and emphasize reading
predicted values from the fitted line, not individual points.

\begin{center}\rule{0.5\linewidth}{0.5pt}\end{center}

\subsection{4.6 Interpreting slope in
context}\label{interpreting-slope-in-context}

Suppose:

\[
y=2.5x+10
\]

where:

\(x\) = hours studied

\(y\) = predicted exam score.

Then \(2.5\) means:

\[
\boxed{\text{for every additional hour studied, predicted score rises by 2.5 points}}.
\]

Not:

\begin{quote}
2.5 extra hours for every point.
\end{quote}

Always track which variable is which.

\begin{center}\rule{0.5\linewidth}{0.5pt}\end{center}

\subsection{4.7 Intercept in context}\label{intercept-in-context}

Suppose:

\[
H=3t+50,
\]

where \(H\) is plant height and \(t\) is days since planting.

At:

\[
t=0,
\]

\[
H=50.
\]

Interpretation:

\[
\boxed{\text{predicted plant height at planting was 50 cm}}.
\]

SAT questions want contextual interpretation, not merely:

\begin{quote}
\(y=50\) when \(x=0\).
\end{quote}

\begin{center}\rule{0.5\linewidth}{0.5pt}\end{center}

\subsection{4.8 Prediction using line of best
fit}\label{prediction-using-line-of-best-fit}

Suppose:

\[
y=.8x+3.
\]

At:

\[
x=10,
\]

\[
y=.8(10)+3=11.
\]

If an actual data point at \(x=10\) happens to be \(12\), the
\textbf{model prediction} is still:

\[
\boxed{11}.
\]

If question says:

\begin{quote}
based on the line of best fit
\end{quote}

use the line.

\begin{center}\rule{0.5\linewidth}{0.5pt}\end{center}

\subsection{4.9 Residual}\label{residual}

Memorize:

\[
\boxed{\text{residual}=\text{actual}-\text{predicted}}.
\]

Think:

\[
\boxed{A-P}.
\]

Example:

actual:

\[
15.
\]

Predicted:

\[
13.
\]

Residual:

\[
15-13=2.
\]

Positive residual.

The actual point is above the line.

Therefore the model \textbf{underestimated}.

\begin{center}\rule{0.5\linewidth}{0.5pt}\end{center}

\subsection{4.10 Residual sign
interpretation}\label{residual-sign-interpretation}

\[
\boxed{\text{positive residual}
\Rightarrow
\text{actual}>\text{predicted}}
\]

so model underestimates.

\[
\boxed{\text{negative residual}
\Rightarrow
\text{actual}<\text{predicted}}
\]

so model overestimates.

\begin{center}\rule{0.5\linewidth}{0.5pt}\end{center}

\subsection{4.11 Residual plots}\label{residual-plots}

A good linear model produces residuals randomly scattered around:

\[
y=0.
\]

No obvious pattern.

If the residual plot is U-shaped:

the linear model missed systematic curvature.

Likely use quadratic/nonlinear model.

Your file specifically flags curvature in the residual plot as evidence
that a linear model is inappropriate.

\begin{center}\rule{0.5\linewidth}{0.5pt}\end{center}

\subsection{4.12 Linear vs exponential}\label{linear-vs-exponential}

This is a classic.

\subsubsection{Linear}\label{linear}

Same amount added each time.

Constant \textbf{difference}.

Example:

\[
100,115,130,145.
\]

Differences:

\[
15,15,15.
\]

Linear.

\begin{center}\rule{0.5\linewidth}{0.5pt}\end{center}

\subsubsection{Exponential}\label{exponential}

Same factor multiplied each time.

Constant \textbf{ratio}.

Example:

\[
3,6,12,24.
\]

Ratios:

\[
2,2,2.
\]

Exponential.

So remember:

\[
\boxed{\text{constant differences}\Rightarrow\text{linear}}
\]

\[
\boxed{\text{constant ratios}\Rightarrow\text{exponential}}.
\]

The source explicitly uses this as the comparison test.

\begin{center}\rule{0.5\linewidth}{0.5pt}\end{center}

\subsection{4.13 Choosing a model}\label{choosing-a-model}

{\def\LTcaptype{none} % do not increment counter
\begin{longtable}[]{@{}ll@{}}
\toprule\noalign{}
Pattern & Likely model \\
\midrule\noalign{}
\endhead
\bottomrule\noalign{}
\endlastfoot
Constant rate/change & Linear \\
One max/min & Quadratic \\
Constant percent/factor & Exponential \\
Rapid rise then leveling & Logarithmic-type \\
\end{longtable}
}

SAT clue words:

``increases by 5 each year'' to linear.

``increases by 5\% each year'' to exponential.

``doubles'' to exponential.

``reaches a maximum'' to quadratic.

\begin{center}\rule{0.5\linewidth}{0.5pt}\end{center}

\subsection{4.14 Interpolation}\label{interpolation}

Interpolation means prediction \textbf{inside} observed data range.

If data ranges:

\[
10\le x\le50,
\]

a prediction at:

\[
x=35
\]

is interpolation.

Usually more reliable.

\begin{center}\rule{0.5\linewidth}{0.5pt}\end{center}

\subsection{4.15 Extrapolation}\label{extrapolation}

Prediction \textbf{outside} observed range.

If data ends at \(x=50\) and you predict \(x=100\):

extrapolation.

The equation may produce a number, but the real-world relationship may
stop behaving that way.

The source explicitly distinguishes interpolation from extrapolation and
notes that extrapolation is less reliable.

\begin{center}\rule{0.5\linewidth}{0.5pt}\end{center}

\section{Part 5: PROBABILITY}\label{part-5-probability}

\subsection{5.1 Basic probability}\label{basic-probability}

\[
\boxed{
P(E)=
\frac{\text{favorable outcomes}}
{\text{total outcomes}}
}
\]

Example:

3 red, 5 blue, 2 green.

Total:

\[
10.
\]

Probability blue:

\[
\frac5{10}
=
\boxed{\frac12}.
\]

\begin{center}\rule{0.5\linewidth}{0.5pt}\end{center}

\subsection{5.2 Complement}\label{complement}

\[
\boxed{P(\text{not }A)=1-P(A)}.
\]

This becomes extremely useful for ``at least one.''

\begin{center}\rule{0.5\linewidth}{0.5pt}\end{center}

\subsection{5.3 ``At least one''}\label{at-least-one}

Usually:

\[
\boxed{
P(\text{at least one})
=
1-P(\text{none})
}
\]

Example:

Roll a die twice.

Probability at least one 6.

Rather than enumerate:

find probability of \textbf{no sixes}.

Each roll:

\[
\frac56.
\]

Both:

\[
\left(\frac56\right)^2
=
\frac{25}{36}.
\]

Therefore:

\[
1-\frac{25}{36}
=
\boxed{\frac{11}{36}}.
\]

This is much safer than adding probabilities.

\begin{center}\rule{0.5\linewidth}{0.5pt}\end{center}

\subsection{5.4 Addition rule}\label{addition-rule}

For events \(A,B\):

\[
\boxed{
P(A\cup B)
=
P(A)+P(B)-P(A\cap B)
}
\]

Why subtract?

Because overlap was counted twice.

\begin{center}\rule{0.5\linewidth}{0.5pt}\end{center}

\subsection{5.5 Mutually exclusive}\label{mutually-exclusive}

Mutually exclusive events cannot both happen.

Therefore:

\[
P(A\cap B)=0.
\]

So:

\[
\boxed{
P(A\cup B)
=
P(A)+P(B)
}
\]

Example:

one die roll.

A = roll a 1.

B = roll a 6.

Cannot both happen on the same roll.

\begin{center}\rule{0.5\linewidth}{0.5pt}\end{center}

\subsection{5.6 Independent events}\label{independent-events}

Independent means event B does not change the probability of A.

Formula:

\[
\boxed{
P(A\cap B)
=
P(A)P(B)
}
\]

Equivalent condition:

\[
\boxed{P(A|B)=P(A)}.
\]

\begin{center}\rule{0.5\linewidth}{0.5pt}\end{center}

\subsection{5.7 With replacement}\label{with-replacement}

Bag has:

3 red,

2 blue.

Draw red, replace it, draw again.

First red:

\[
\frac35.
\]

Since replaced, second red:

\[
\frac35.
\]

Both:

\[
\frac35\cdot\frac35
=
\boxed{\frac9{25}}.
\]

\begin{center}\rule{0.5\linewidth}{0.5pt}\end{center}

\subsection{5.8 Without replacement}\label{without-replacement}

Same bag.

Draw two without replacement.

First red:

\[
\frac35.
\]

Now only 2 red remain out of 4 balls.

Second red:

\[
\frac24.
\]

Therefore:

\[
\frac35\cdot\frac24
=
\frac6{20}
=
\boxed{\frac3{10}}.
\]

Without replacement creates dependence because probabilities change.

\begin{center}\rule{0.5\linewidth}{0.5pt}\end{center}

\subsection{5.9 Conditional probability}\label{conditional-probability}

This is one of the most important SAT probability topics.

Formula:

\[
\boxed{
P(A|B)
=
\frac{P(A\cap B)}
{P(B)}
}
\]

But for SAT tables, there is an easier rule:

\begin{quote}
\textbf{The denominator is the group after ``given that.''}
\end{quote}

Your source explicitly teaches this denominator rule.

\begin{center}\rule{0.5\linewidth}{0.5pt}\end{center}

\subsection{5.10 Two-way table example}\label{two-way-table-example}

{\def\LTcaptype{none} % do not increment counter
\begin{longtable}[]{@{}lrrr@{}}
\toprule\noalign{}
& Passed & Failed & Total \\
\midrule\noalign{}
\endhead
\bottomrule\noalign{}
\endlastfoot
Studied & 100 & 20 & 120 \\
Didn't study & 30 & 50 & 80 \\
Total & 130 & 70 & 200 \\
\end{longtable}
}

Question:

\begin{quote}
Given that a student studied, what is the probability they passed?
\end{quote}

The condition is:

\[
\boxed{\text{studied}}.
\]

Therefore denominator:

\[
120.
\]

Studied AND passed:

\[
100.
\]

So:

\[
P(\text{passed}|\text{studied})
=
\frac{100}{120}
=
\boxed{\frac56}.
\]

Do NOT use:

\[
\frac{100}{200}.
\]

That ignores the condition.

\begin{center}\rule{0.5\linewidth}{0.5pt}\end{center}

\subsection{5.11 Reverse condition trap}\label{reverse-condition-trap}

From the same table:

\[
P(\text{studied}|\text{passed})
=
\frac{100}{130}.
\]

This is different from:

\[
P(\text{passed}|\text{studied})
=
\frac{100}{120}.
\]

Conditional probability is not generally reversible.

\[
\boxed{P(A|B)\neq P(B|A)}
\]

in general.

\begin{center}\rule{0.5\linewidth}{0.5pt}\end{center}

\subsection{5.12 Independence from a
table}\label{independence-from-a-table}

Using that table:

\[
P(\text{passed})
=
\frac{130}{200}=0.65.
\]

But:

\[
P(\text{passed}|\text{studied})
=
\frac{100}{120}
\approx0.833.
\]

Since:

\[
0.833\neq0.65,
\]

studying and passing are \textbf{not independent}.

\begin{center}\rule{0.5\linewidth}{0.5pt}\end{center}

\subsection{5.13 Reverse conditional
probability}\label{reverse-conditional-probability}

Suppose:

80\% of the 120 students who studied passed.

Number who passed and studied:

\[
0.80(120)=96.
\]

\[
\boxed{96}.
\]

This appears in your notes as the ``reverse probability'' form.

\begin{center}\rule{0.5\linewidth}{0.5pt}\end{center}

\subsection{5.14 Geometric probability}\label{geometric-probability}

When a point is uniformly random over an area:

\[
\boxed{
P(E)
=
\frac{\text{favorable area}}
{\text{total area}}
}
\]

Example:

square side:

\[
20.
\]

Circle radius:

\[
5.
\]

Square area:

\[
400.
\]

Circle:

\[
25\pi.
\]

Probability:

\[
\frac{25\pi}{400}
=
\boxed{\frac{\pi}{16}}.
\]

Do not compare lengths.

Use \textbf{areas}.

\begin{center}\rule{0.5\linewidth}{0.5pt}\end{center}

\section{Part 6: STATISTICAL
INFERENCE}\label{part-6-statistical-inference}

This section is less computational and more about exact wording.

\subsection{6.1 Population vs sample}\label{population-vs-sample}

Population:

\[
\boxed{\text{entire group we care about}}.
\]

Sample:

\[
\boxed{\text{subset actually studied}}.
\]

Example:

A school has 1,500 students.

Researchers survey 200.

Population:

\[
1500.
\]

Sample:

\[
200.
\]

\begin{center}\rule{0.5\linewidth}{0.5pt}\end{center}

\subsection{6.2 Parameter vs statistic}\label{parameter-vs-statistic}

A \textbf{parameter} describes a population.

A \textbf{statistic} describes a sample.

Example:

42\% of surveyed voters support something.

That:

\[
42\%
\]

is a sample statistic.

We use it to estimate the true population parameter.

The source frames this exactly as ``sample statistic to estimate of
population parameter.''

\begin{center}\rule{0.5\linewidth}{0.5pt}\end{center}

\subsection{6.3 Estimating a population
total}\label{estimating-a-population-total}

Suppose:

60 of 200 sampled students own cars.

Sample proportion:

\[
\frac{60}{200}=0.30.
\]

School has 1,500 students.

Estimate:

\[
0.30(1500)=450.
\]

So approximately:

\[
\boxed{450}.
\]

Notice the word:

\textbf{approximately}.

Sampling produces estimates.

\begin{center}\rule{0.5\linewidth}{0.5pt}\end{center}

\subsection{6.4 Margin of error}\label{margin-of-error}

Suppose a poll says:

\[
42\%
\]

with margin of error:

\[
\pm4\%.
\]

Then plausible population range:

\[
42-4=38
\]

to:

\[
42+4=46.
\]

Therefore:

\[
\boxed{38\%-46\%}.
\]

General:

\[
\boxed{
\text{estimate}\pm\text{MoE}
}
\]

\begin{center}\rule{0.5\linewidth}{0.5pt}\end{center}

\subsection{6.5 What margin of error DOES NOT
mean}\label{what-margin-of-error-does-not-mean}

If the sample found exactly 42\%, you should not say:

\begin{quote}
Between 38\% and 46\% of the surveyed people supported it.
\end{quote}

The surveyed proportion is already known.

The uncertainty concerns the \textbf{population parameter}.

Correct idea:

\begin{quote}
The population proportion is plausibly between 38\% and 46\%.
\end{quote}

The notes explicitly warn that the interval concerns the parameter, not
individual observations or the already-observed sample.

\begin{center}\rule{0.5\linewidth}{0.5pt}\end{center}

\subsection{6.6 Sample size and margin of
error}\label{sample-size-and-margin-of-error}

Holding other things comparable:

larger sample:

\[
\boxed{\text{smaller margin of error}}
\]

which means:

\[
\boxed{\text{more precision}}.
\]

Smaller sample:

larger margin of error.

But:

\subsection{HUGE distinction}\label{huge-distinction}

A very large \textbf{biased sample is still biased}.

Sampling 1 million people from the wrong group does not magically make
the result representative.

\begin{center}\rule{0.5\linewidth}{0.5pt}\end{center}

\subsection{6.7 Confidence intervals}\label{confidence-intervals}

Suppose:

\[
95\%\text{ CI}=(122,138)
\]

for a population mean.

SAT interpretation:

\[
\boxed{\text{the true population mean is plausibly between 122 and 138}}.
\]

Do NOT say:

\begin{quote}
95\% of people have values between 122 and 138.
\end{quote}

The interval concerns the \textbf{mean/parameter}, not individual
observations.

\begin{center}\rule{0.5\linewidth}{0.5pt}\end{center}

\subsection{6.8 What ``95\% confidence''
means}\label{what-95-confidence-means}

Conceptually:

If the same sampling procedure were repeated many times, about 95\% of
intervals constructed this way would capture the true population
parameter.

For SAT purposes, the most important thing is:

\[
\boxed{\text{confidence interval is about population parameter}}
\]

not:

individual values.

\begin{center}\rule{0.5\linewidth}{0.5pt}\end{center}

\subsection{6.9 Random sample}\label{random-sample}

Random selection allows you to \textbf{generalize}.

Suppose you randomly sample registered voters in Ohio.

You can generalize to:

\[
\boxed{\text{registered voters in Ohio}}.
\]

Not automatically:

all Americans.

Not automatically:

all Ohio residents.

The population is the group from which the sample was selected.

Your file explicitly distinguishes random selection as enabling
generalization.

\begin{center}\rule{0.5\linewidth}{0.5pt}\end{center}

\subsection{6.10 Convenience sample}\label{convenience-sample}

Example:

Want to learn how often Americans exercise.

Survey people entering a gym.

That is biased.

Why?

People at a gym are likely more exercise-oriented than average
Americans.

Even if you survey 10,000 gym members, it is still not representative.

\begin{center}\rule{0.5\linewidth}{0.5pt}\end{center}

\subsection{6.11 Voluntary-response bias}\label{voluntary-response-bias}

You post:

\begin{quote}
Vote online if you want to tell us what you think!
\end{quote}

People with stronger opinions may be more likely to respond.

Therefore voluntary response is often biased.

\begin{center}\rule{0.5\linewidth}{0.5pt}\end{center}

\subsection{6.12 Nonresponse bias}\label{nonresponse-bias}

You randomly choose people, but most do not answer.

If responders differ systematically from nonresponders, your results can
be biased.

\begin{center}\rule{0.5\linewidth}{0.5pt}\end{center}

\subsection{6.13 Undercoverage}\label{undercoverage}

Part of the target population has little/no chance of appearing in the
sample.

Example:

A survey conducted only by landline may underrepresent people without
landlines.

\begin{center}\rule{0.5\linewidth}{0.5pt}\end{center}

\subsection{6.14 Leading/question wording
bias}\label{leadingquestion-wording-bias}

Question:

\begin{quote}
Do you support this sensible and effective new policy?
\end{quote}

Words like ``sensible and effective'' push the respondent.

That can bias responses.

\begin{center}\rule{0.5\linewidth}{0.5pt}\end{center}

\subsection{6.15 Social desirability
bias}\label{social-desirability-bias}

People may give answers they think are socially acceptable rather than
fully truthful.

Example:

``How often do you recycle?''

Respondents may overreport.

\begin{center}\rule{0.5\linewidth}{0.5pt}\end{center}

\section{Part 7: STUDY DESIGN AND
CAUSATION}\label{part-7-study-design-and-causation}

This is probably the easiest conceptual area once you learn one
distinction.

\subsection{7.1 The two words you must never
confuse}\label{the-two-words-you-must-never-confuse}

\subsection{Random selection}\label{random-selection}

Answers:

\begin{quote}
Who gets into the study?
\end{quote}

Enables:

\[
\boxed{\text{generalization}}
\]

\begin{center}\rule{0.5\linewidth}{0.5pt}\end{center}

\subsection{Random assignment}\label{random-assignment}

Answers:

\begin{quote}
Once people are in the study, who gets which treatment?
\end{quote}

Enables:

\[
\boxed{\text{causal inference}}
\]

This is the most important study-design distinction in the domain. Your
source's formula sheet summarizes it as random selection to generalize,
random assignment to causation.

\begin{center}\rule{0.5\linewidth}{0.5pt}\end{center}

\subsection{7.2 Four possible
situations}\label{four-possible-situations}

{\def\LTcaptype{none} % do not increment counter
\begin{longtable}[]{@{}lll@{}}
\toprule\noalign{}
Random selection? & Random assignment? & What can you do? \\
\midrule\noalign{}
\endhead
\bottomrule\noalign{}
\endlastfoot
Yes & Yes & Generalize + causal claim \\
Yes & No & Generalize, association only \\
No & Yes & Causal claim for studied subjects, limited generalization \\
No & No & Association only, limited generalization \\
\end{longtable}
}

Memorize this table.

\begin{center}\rule{0.5\linewidth}{0.5pt}\end{center}

\subsection{7.3 Observational study}\label{observational-study}

Researchers merely observe what people already do.

Example:

Researchers compare people who drink coffee versus those who don't.

Nobody is assigned coffee.

Therefore:

\[
\boxed{\text{association only}}
\]

You cannot automatically conclude:

coffee causes the observed outcome.

\begin{center}\rule{0.5\linewidth}{0.5pt}\end{center}

\subsection{7.4 Experiment}\label{experiment}

Researchers impose a treatment.

For example:

randomly assign patients to:

drug,

placebo.

If outcomes differ significantly, random assignment supports a causal
interpretation.

\begin{center}\rule{0.5\linewidth}{0.5pt}\end{center}

\subsection{7.5 Correlation does not establish
causation}\label{correlation-does-not-establish-causation}

Suppose people who exercise more have lower blood pressure.

Maybe exercise lowers blood pressure.

But observationally there could be other explanations:

diet,

age,

income,

health awareness,

medication,

etc.

So observational result:

\[
\boxed{\text{exercise is associated with blood pressure}}
\]

not necessarily:

\[
\boxed{\text{exercise caused the difference}}.
\]

\begin{center}\rule{0.5\linewidth}{0.5pt}\end{center}

\subsection{7.6 Confounding variable}\label{confounding-variable}

A confounder is a third variable related to both:

the explanatory/treatment variable,

and the outcome.

Classic example:

Ice cream sales increase when drowning rates increase.

Did ice cream cause drowning?

No.

Temperature is a confounder:

hot weather to more ice cream.

hot weather to more swimming.

more swimming to potentially more drownings.

So temperature helps explain both.

\begin{center}\rule{0.5\linewidth}{0.5pt}\end{center}

\subsection{7.7 Another confounding
example}\label{another-confounding-example}

Students taking SAT-prep courses have higher scores.

Possible confounder:

\[
\boxed{\text{motivation}}.
\]

More motivated students may be:

more likely to enroll in prep,

and more likely to study independently and score higher.

So an observational association does not prove the course caused the
entire difference.

\begin{center}\rule{0.5\linewidth}{0.5pt}\end{center}

\subsection{7.8 Control group}\label{control-group}

A control group gives a baseline.

Example:

Treatment group gets new medicine.

Control group gets existing treatment or placebo.

Then compare outcomes.

\begin{center}\rule{0.5\linewidth}{0.5pt}\end{center}

\subsection{7.9 Placebo}\label{placebo}

A fake/inactive treatment designed to resemble the real one.

Why?

People may improve simply because they believe they are being treated.

That is the placebo effect.

\begin{center}\rule{0.5\linewidth}{0.5pt}\end{center}

\subsection{7.10 Blinding}\label{blinding}

Single-blind:

participants do not know their assigned group.

Double-blind:

participants and relevant researchers/evaluators do not know.

Purpose:

reduce expectation and measurement bias.

\begin{center}\rule{0.5\linewidth}{0.5pt}\end{center}

\subsection{7.11 Blocking}\label{blocking}

Researchers first group subjects by an important characteristic:

age,

sex,

baseline severity,

etc.

Then randomize within groups.

Purpose:

balance known sources of variation.

\begin{center}\rule{0.5\linewidth}{0.5pt}\end{center}

\section{Part 8: HIGH-YIELD SAT
WORDING}\label{part-8-high-yield-sat-wording}

These words should trigger methods automatically.

{\def\LTcaptype{none} % do not increment counter
\begin{longtable}[]{@{}ll@{}}
\toprule\noalign{}
SAT wording & Think \\
\midrule\noalign{}
\endhead
\bottomrule\noalign{}
\endlastfoot
``per'' & rate \\
``varies directly'' & \(y=kx\) \\
``varies inversely'' & \(y=k/x\) \\
``increase by \(p\%\)'' & \(\times(1+p/100)\) \\
``decrease by \(p\%\)'' & \(\times(1-p/100)\) \\
``after a discount'' & reverse percent may require division \\
``percent change'' & divide change by old \\
``percentage points'' & subtract percentages \\
``mean'' & sum / count \\
``middle'' & median after sorting \\
``spread'' & range/IQR/SD \\
``resistant to outliers'' & median/IQR \\
``line of best fit'' & use model, not data point \\
``predicted'' & model value \\
``residual'' & actual − predicted \\
``given that'' & conditional probability \\
``among those who'' & conditional denominator \\
``at least one'' & \(1-P(\text{none})\) \\
``randomly selected'' & generalization \\
``randomly assigned'' & causation \\
``observational'' & association only \\
``margin of error'' & estimate ± MoE \\
``plausibly'' & good inference language \\
``exactly'' & suspicious for sample inference \\
\end{longtable}
}

\begin{center}\rule{0.5\linewidth}{0.5pt}\end{center}

\section{Part 9: DESMOS FOR B3}\label{part-9-desmos-for-b3}

You should absolutely use Desmos for some Data Analysis problems.

\subsection{9.1 Mean/median/list
calculations}\label{meanmedianlist-calculations}

Enter:

\begin{Shaded}
\begin{Highlighting}[]
\NormalTok{L=[12,7,9,14,3]}
\end{Highlighting}
\end{Shaded}

Then:

\begin{Shaded}
\begin{Highlighting}[]
\NormalTok{mean(L)}
\end{Highlighting}
\end{Shaded}

or:

\begin{Shaded}
\begin{Highlighting}[]
\NormalTok{median(L)}
\end{Highlighting}
\end{Shaded}

\begin{center}\rule{0.5\linewidth}{0.5pt}\end{center}

\subsection{9.2 Unit conversions}\label{unit-conversions}

For:

60 mph to ft/s:

\begin{Shaded}
\begin{Highlighting}[]
\NormalTok{60*5280/3600}
\end{Highlighting}
\end{Shaded}

Answer:

\[
88.
\]

\begin{center}\rule{0.5\linewidth}{0.5pt}\end{center}

\subsection{9.3 Percentage chains}\label{percentage-chains}

30\% markup then 15\% discount:

\begin{Shaded}
\begin{Highlighting}[]
\NormalTok{150*1.30*0.85}
\end{Highlighting}
\end{Shaded}

\begin{center}\rule{0.5\linewidth}{0.5pt}\end{center}

\subsection{9.4 Reverse percentage}\label{reverse-percentage-1}

\$45 after 25\% discount:

\begin{Shaded}
\begin{Highlighting}[]
\NormalTok{45/0.75}
\end{Highlighting}
\end{Shaded}

\begin{center}\rule{0.5\linewidth}{0.5pt}\end{center}

\subsection{9.5 Linear regression}\label{linear-regression}

Create table \(x_1,y_1\).

Then:

\begin{Shaded}
\begin{Highlighting}[]
\NormalTok{y\_1 \textasciitilde{} m x\_1 + b}
\end{Highlighting}
\end{Shaded}

Desmos gives \(m,b\).

\begin{center}\rule{0.5\linewidth}{0.5pt}\end{center}

\subsection{9.6 Quadratic regression}\label{quadratic-regression}

\begin{Shaded}
\begin{Highlighting}[]
\NormalTok{y\_1 \textasciitilde{} a x\_1\^{}2 + b x\_1 + c}
\end{Highlighting}
\end{Shaded}

\begin{center}\rule{0.5\linewidth}{0.5pt}\end{center}

\subsection{9.7 Exponential regression}\label{exponential-regression}

\begin{Shaded}
\begin{Highlighting}[]
\NormalTok{y\_1 \textasciitilde{} a*b\^{}(x\_1)}
\end{Highlighting}
\end{Shaded}

Your source recommends regression and comparing fits as one of the main
calculator workflows for B3.

\begin{center}\rule{0.5\linewidth}{0.5pt}\end{center}

\section{Part 10: THE BIGGEST B3
TRAPS}\label{part-10-the-biggest-b3-traps}

Here are the ones I want you actively checking for on test day.

{\def\LTcaptype{none} % do not increment counter
\begin{longtable}[]{@{}ll@{}}
\toprule\noalign{}
Trap & Correct rule \\
\midrule\noalign{}
\endhead
\bottomrule\noalign{}
\endlastfoot
Percent change divided by new value & Divide by \textbf{old} \\
Successive percents added & Multiply multipliers \\
20\% rise then 20\% fall = same & No, multiplier \(1.2(.8)=.96\) \\
Area conversion uses linear factor & Square conversion factor \\
Volume conversion uses linear factor & Cube conversion factor \\
Average speeds & Total distance / total time \\
Weighted groups & Weight by group size \\
Median before sorting & Sort first \\
Same range = same SD & False \\
Outlier affects mean and median equally & False; mean more sensitive \\
Right skew means median \textgreater{} mean & Usually mean
\textgreater{} median \\
Using data point when asked model prediction & Use fitted line/model \\
Residual = predicted − actual & Actual − predicted \\
Positive residual = overestimate & Positive means underestimate \\
Correlation proves causation & False \\
``Given B'' uses grand total & Use total of B \\
At least one to add probabilities & Usually use complement \\
Without replacement probabilities stay same & They change \\
Large sample fixes selection bias & False \\
Random sample proves causation & False \\
Random assignment allows generalization & Not by itself \\
CI describes individual observations & It describes parameter \\
Sample statistic is exact population truth & It is an estimate \\
\end{longtable}
}

\begin{center}\rule{0.5\linewidth}{0.5pt}\end{center}

\section{Part 11: COMPLETE FORMULA
SHEET}\label{part-11-complete-formula-sheet}

\subsection{Ratios/proportions}\label{ratiosproportions}

\[
\boxed{\frac ab=\frac cd\Rightarrow ad=bc}
\]

\subsection{Direct variation}\label{direct-variation-1}

\[
\boxed{y=kx}
\]

\subsection{Inverse variation}\label{inverse-variation-1}

\[
\boxed{y=\frac{k}{x}}
\]

or:

\[
\boxed{xy=k}
\]

\subsection{Density}\label{density-1}

\[
\boxed{D=\frac MV}
\]

\subsection{Speed}\label{speed}

\[
\boxed{v=\frac dt}
\]

\subsection{Average speed}\label{average-speed-1}

\[
\boxed{
v_{\text{avg}}
=
\frac{\text{total distance}}
{\text{total time}}
}
\]

\subsection{Weighted average}\label{weighted-average-1}

\[
\boxed{
\bar x=
\frac{\sum w_ix_i}{\sum w_i}
}
\]

\begin{center}\rule{0.5\linewidth}{0.5pt}\end{center}

\subsection{Percent of}\label{percent-of}

\[
\boxed{p\%\text{ of }x=\frac p{100}x}
\]

\subsection{Percent change}\label{percent-change-1}

\[
\boxed{
\frac{\text{new}-\text{old}}
{\text{old}}\times100\%
}
\]

\subsection{Increase multiplier}\label{increase-multiplier}

\[
\boxed{1+\frac p{100}}
\]

\subsection{Decrease multiplier}\label{decrease-multiplier}

\[
\boxed{1-\frac p{100}}
\]

\subsection{Reverse percent}\label{reverse-percent}

\[
\boxed{
\text{original}
=
\frac{\text{final}}
{\text{multiplier}}
}
\]

\subsection{Compound growth}\label{compound-growth-1}

\[
\boxed{A=P(1+r)^n}
\]

\subsection{Simple interest}\label{simple-interest-1}

\[
\boxed{I=Prt}
\]

\subsection{Percent error}\label{percent-error-1}

\[
\boxed{
\frac{|\text{estimate}-\text{actual}|}
{\text{actual}}\times100\%
}
\]

\begin{center}\rule{0.5\linewidth}{0.5pt}\end{center}

\subsection{Mean}\label{mean-1}

\[
\boxed{\bar x=\frac{\sum x}{n}}
\]

\subsection{Sum from mean}\label{sum-from-mean}

\[
\boxed{\sum x=n\bar x}
\]

\subsection{Range}\label{range-1}

\[
\boxed{\max-\min}
\]

\subsection{IQR}\label{iqr}

\[
\boxed{Q_3-Q_1}
\]

\subsection{Outlier fences}\label{outlier-fences}

\[
\boxed{Q_1-1.5(IQR)}
\]

and:

\[
\boxed{Q_3+1.5(IQR)}
\]

\begin{center}\rule{0.5\linewidth}{0.5pt}\end{center}

\subsection{Linear model}\label{linear-model}

\[
\boxed{y=mx+b}
\]

\subsection{Residual}\label{residual-1}

\[
\boxed{\text{actual}-\text{predicted}}
\]

\subsection{Exponential model}\label{exponential-model}

\[
\boxed{y=ab^x}
\]

\begin{center}\rule{0.5\linewidth}{0.5pt}\end{center}

\subsection{Basic probability}\label{basic-probability-1}

\[
\boxed{
P(E)=
\frac{\text{favorable}}
{\text{total}}
}
\]

\subsection{Complement}\label{complement-1}

\[
\boxed{P(A^c)=1-P(A)}
\]

\subsection{Addition}\label{addition}

\[
\boxed{
P(A\cup B)
=
P(A)+P(B)-P(A\cap B)
}
\]

\subsection{Conditional}\label{conditional}

\[
\boxed{
P(A|B)
=
\frac{P(A\cap B)}
{P(B)}
}
\]

\subsection{Independent}\label{independent}

\[
\boxed{
P(A\cap B)=P(A)P(B)
}
\]

\subsection{At least one}\label{at-least-one-1}

\[
\boxed{
P(\ge1)=1-P(0)
}
\]

\subsection{Geometric probability}\label{geometric-probability-1}

\[
\boxed{
P(E)=
\frac{\text{favorable area}}
{\text{total area}}
}
\]

\begin{center}\rule{0.5\linewidth}{0.5pt}\end{center}

\subsection{Margin of error}\label{margin-of-error-1}

\[
\boxed{
\text{interval}
=
\text{statistic}\pm\text{MoE}
}
\]

\begin{center}\rule{0.5\linewidth}{0.5pt}\end{center}

\section{Part 12: FINAL SAT B3 DECISION
TREE}\label{part-12-final-sat-b3-decision-tree}

This is the final checklist I would study right before the test:

\begin{enumerate}
\def\labelenumi{\arabic{enumi}.}
\tightlist
\item
  \textbf{Ratio?} Decide part-to-part or part-to-whole.
\item
  \textbf{Proportion?} Keep matching units in corresponding positions.
\item
  \textbf{Direct variation?} Use \(y=kx\).
\item
  \textbf{Inverse variation?} Use \(xy=k\).
\item
  \textbf{Unit conversion?} Make units cancel.
\item
  \textbf{Area/volume conversion?} Square/cube the factor.
\item
  \textbf{Average speed?} Total distance / total time.
\item
  \textbf{Weighted groups?} Use weighted average.
\item
  \textbf{Percent of?} Multiply by decimal.
\item
  \textbf{Percent change?} \((new-old)/old\).
\item
  \textbf{Increase/decrease?} Use multiplier.
\item
  \textbf{Several percent changes?} Multiply multipliers.
\item
  \textbf{Original before discount?} Divide by multiplier.
\item
  \textbf{Percentage points?} Subtract percentage values.
\item
  \textbf{Mean?} Sum/count.
\item
  \textbf{Missing value with mean?} Mean × count = total.
\item
  \textbf{Median?} Sort first.
\item
  \textbf{IQR?} \(Q_3-Q_1\).
\item
  \textbf{Outlier?} 1.5×IQR rule.
\item
  \textbf{Outlier present?} Mean and SD are sensitive; median/IQR
  resistant.
\item
  \textbf{Right skew?} Mean usually above median.
\item
  \textbf{Left skew?} Mean usually below median.
\item
  \textbf{Compare SD?} Look at distances from center, not just range.
\item
  \textbf{Add constant to all data?} Center moves; spread stays.
\item
  \textbf{Multiply all data?} Center and spread scale.
\item
  \textbf{Line of best fit?} Use the line, not a nearby point.
\item
  \textbf{Slope?} Predicted \(y\)-change per +1 \(x\).
\item
  \textbf{Residual?} Actual − predicted.
\item
  \textbf{Residual positive?} Model underestimated.
\item
  \textbf{Curved residual plot?} Linear model inappropriate.
\item
  \textbf{Constant differences?} Linear.
\item
  \textbf{Constant ratios?} Exponential.
\item
  \textbf{Prediction outside data range?} Extrapolation.
\item
  \textbf{Probability?} Identify favorable and total outcomes.
\item
  \textbf{``Given that''/``among''?} Conditional probability; change
  denominator.
\item
  \textbf{Without replacement?} Update probabilities.
\item
  \textbf{At least one?} Complement method.
\item
  \textbf{Mutually exclusive?} Intersection = 0.
\item
  \textbf{Independent?} \(P(A|B)=P(A)\).
\item
  \textbf{Area probability?} Favorable area / total area.
\item
  \textbf{Sample result?} It estimates a population parameter.
\item
  \textbf{Margin of error?} Statistic ± MoE.
\item
  \textbf{Larger random sample?} Usually smaller MoE.
\item
  \textbf{Huge but biased sample?} Still biased.
\item
  \textbf{Randomly selected?} Can generalize to sampled population.
\item
  \textbf{Randomly assigned?} Can support causation.
\item
  \textbf{Observational study?} Association only.
\item
  \textbf{Confounder?} Related to both explanatory variable and outcome.
\item
  \textbf{Volunteers/convenience sample?} Limited generalizability.
\item
  \textbf{Answer says ``exactly'' about population from sample?} Usually
  reject it.
\end{enumerate}

If those 50 reactions become automatic, \textbf{B3 stops being a
collection of random word problems}. You start recognizing a relatively
small number of repeatable structures, which is exactly what this domain
is designed around.

\end{document}
